\documentclass{article}
\usepackage{amsmath, amssymb, amsthm}
\usepackage{hyperref}
\usepackage{geometry}
\geometry{margin=1in}

\title{Erd\H{o}s Problem \#349}
\date{Last updated: 2025-08-31}
\author{Source: erdosproblems.com}

\begin{document}
\maketitle

\noindent\textbf{Status:} OPEN \quad \textbf{No prize}

\noindent\textbf{Tags:} number theory, complete sequences

\medskip
\noindent\textbf{Problem Statement:}

For what values of $t,\alpha \in (0,\infty)$ is the sequence $\lfloor t\alpha^n\rfloor$ complete (that is, all sufficiently large integers are the sum of distinct integers of the form $\lfloor t\alpha^n\rfloor$)?

\bigskip
\noindent\textbf{Remarks:}

Even in the range $t\in (0,1)$ and $\alpha\in (1,2)$ the behaviour is surprisingly complex. For example, Graham \cite{Gr64e} has shown that for any $k$ there exists some $t_k\in (0,1)$ such that the set of $\alpha$ such that the sequence is complete consists of at least $k$ disjoint line segments. It seems likely that the sequence is complete for all $t>0$ and all $1<\alpha < \frac{1+\sqrt{5}}{2}$. Proving this seems very difficult, since we do not even know whether $\lfloor (3/2)^n\rfloor$ is odd or even infinitely often.

\bigskip
\noindent\textbf{References:}

[Gr64e] Graham, R. L., On a conjecture of Erd\H{o}s in additive number theory. Acta Arith. (1964/65), 63-70.

\bigskip
\noindent\small{Source: \url{https://www.erdosproblems.com/349}}
\end{document}
